\begingroup
\renewcommand{\arraystretch}{1.45}
\Needspace{37\baselineskip}
\InterpolationArticle{10}
Es sei \(X\) eine Funktion des unbestimmten Bogens \(x\) von der Form
\[
\begin{aligned}
&\alpha+\alpha'\cos x+\alpha''\cos 2x+\alpha'''\cos 3x+\text{usw.}\\
&\quad+\beta'\sin x+\beta''\sin 2x+\beta'''\sin 3x+\text{usw.}
\end{aligned}
\]
die nicht bis ins Unendliche fortgesetzt wird, sondern mit \(\cos mx\) und \(\sin mx\) abbricht, sodass
die Anzahl der Koeffizienten (Unbekannten) \(2m+1\) beträgt. Für ebenso viele verschiedene Werte von \(x\), nämlich \(a,b,c,d\) usw., seien die entsprechenden Werte der Funktion \(X\) gegeben,
nämlich \(A,B,C,D\mathinner{.\,.\,.}\) (Werte von \(x\), deren Differenz ein ganzer Umfang oder ein Vielfaches davon ist, können hier offenbar nicht als verschieden gelten). Aus
diesen Daten wird eine Formel für den Wert \(T\) gesucht, den die Funktion \(X\) für einen beliebigen anderen
Wert von \(x\), nämlich \(t\), annimmt. Es gelten somit \(2m+2\) Gleichungen
\SourceLeaf{289}
\[
\begin{aligned}
A&=\alpha+\alpha'\cos a+\alpha''\cos 2a+\alpha'''\cos 3a+\text{usw.}\\
 &\quad+\beta'\sin a+\beta''\sin 2a+\beta'''\sin 3a+\text{usw.}\\
B&=\alpha+\alpha'\cos b+\alpha''\cos 2b+\alpha'''\cos 3b+\text{usw.}\\
 &\quad+\beta'\sin b+\beta''\sin 2b+\beta'''\sin 3b+\text{usw.}\\
C&=\alpha+\alpha'\cos c+\alpha''\cos 2c+\alpha'''\cos 3c+\text{usw.}\\
 &\quad+\beta'\sin c+\beta''\sin 2c+\beta'''\sin 3c+\text{usw.}\\
D&=\alpha+\alpha'\cos d+\alpha''\cos 2d+\alpha'''\cos 3d+\text{usw.}\\
 &\quad+\beta'\sin d+\beta''\sin 2d+\beta'''\sin 3d+\text{usw.}\\
 &\text{usw.}\\
T&=\alpha+\alpha'\cos t+\alpha''\cos 2t+\alpha'''\cos 3t+\text{usw.}\\
 &\quad+\beta'\sin t+\beta''\sin 2t+\beta'''\sin 3t+\text{usw.}
\end{aligned}
\]
Man multipliziere diese Gleichungen der Reihe nach mit
\[
\begin{gathered}
\frac{1}{\sin\tfrac12(a-b)\sin\tfrac12(a-c)\sin\tfrac12(a-d)\mathinner{.\,.\,.\,.}\sin\tfrac12(a-t)}\\[6pt]
\frac{1}{\sin\tfrac12(b-a)\sin\tfrac12(b-c)\sin\tfrac12(b-d)\mathinner{.\,.\,.\,.}\sin\tfrac12(b-t)}\\[6pt]
\frac{1}{\sin\tfrac12(c-a)\sin\tfrac12(c-b)\sin\tfrac12(c-d)\mathinner{.\,.\,.\,.}\sin\tfrac12(c-t)}\\[6pt]
\frac{1}{\sin\tfrac12(d-a)\sin\tfrac12(d-b)\sin\tfrac12(d-c)\mathinner{.\,.\,.\,.}\sin\tfrac12(d-t)}\\
\text{usw.}\\[6pt]
\frac{1}{\sin\tfrac12(t-a)\sin\tfrac12(t-b)\sin\tfrac12(t-c)\sin\tfrac12(t-d)\mathinner{.\,.\,.}}
\end{gathered}
\]
und addiere die Produkte, sodass sich ergibt
\[
\begin{aligned}
&\frac{A}{\sin\tfrac12(a-b)\sin\tfrac12(a-c)\sin\tfrac12(a-d)\mathinner{.\,.\,.\,.}\sin\tfrac12(a-t)}\\[6pt]
+{}&\frac{B}{\sin\tfrac12(b-a)\sin\tfrac12(b-c)\sin\tfrac12(b-d)\mathinner{.\,.\,.\,.}\sin\tfrac12(b-t)}\\[6pt]
+{}&\frac{C}{\sin\tfrac12(c-a)\sin\tfrac12(c-b)\sin\tfrac12(c-d)\mathinner{.\,.\,.\,.}\sin\tfrac12(c-t)}\\[6pt]
+{}&\frac{D}{\sin\tfrac12(d-a)\sin\tfrac12(d-b)\sin\tfrac12(d-c)\mathinner{.\,.\,.\,.}\sin\tfrac12(d-t)}\\
+{}&\text{usw.}\\[6pt]
+{}&\frac{T}{\sin\tfrac12(t-a)\sin\tfrac12(t-b)\sin\tfrac12(t-c)\sin\tfrac12(t-d)\mathinner{.\,.\,.\,.}}=W
\end{aligned}
\]
\SourceLeaf{290}
Dann folgt aus Artikel 2, wo $m$ dasselbe war, was hier $2m+2$ ist, sodass der zweite Fall vorliegt, $W=0$; durch Multiplikation mit
\[
\sin\tfrac12(t-a)\sin\tfrac12(t-b)\sin\tfrac12(t-c)\sin\tfrac12(t-d)\mathinner{.\,.\,.}
\]
ergibt sich daher
\[
\begin{aligned}
T={}&\frac{\sin\tfrac12(t-b)\sin\tfrac12(t-c)\sin\tfrac12(t-d)\mathinner{.\,.\,.\,.}}
{\sin\tfrac12(a-b)\sin\tfrac12(a-c)\sin\tfrac12(a-d)\mathinner{.\,.\,.\,.}}A\\[5pt]
&+\frac{\sin\tfrac12(t-a)\sin\tfrac12(t-c)\sin\tfrac12(t-d)\mathinner{.\,.\,.\,.}}
{\sin\tfrac12(b-a)\sin\tfrac12(b-c)\sin\tfrac12(b-d)\mathinner{.\,.\,.\,.}}B\\[5pt]
&+\frac{\sin\tfrac12(t-a)\sin\tfrac12(t-b)\sin\tfrac12(t-d)\mathinner{.\,.\,.\,.}}
{\sin\tfrac12(c-a)\sin\tfrac12(c-b)\sin\tfrac12(c-d)\mathinner{.\,.\,.\,.}}C\\[5pt]
&+\frac{\sin\tfrac12(t-a)\sin\tfrac12(t-b)\sin\tfrac12(t-c)\mathinner{.\,.\,.\,.}}
{\sin\tfrac12(d-a)\sin\tfrac12(d-b)\sin\tfrac12(d-c)\mathinner{.\,.\,.\,.}}D\\[5pt]
&+\text{usw.}
\end{aligned}
\]
\Needspace{16\baselineskip}
Da diese Formel ohne Einschränkung für jeden beliebigen Wert von $t$ gilt, ist offenbar, dass nach Entwicklung der Sinusprodukte in den Zählern in Kosinus und Sinus von Vielfachen
der Bögen das Ergebnis und
\[
\begin{aligned}
&\alpha+\alpha'\cos t+\alpha''\cos2t+\alpha'''\cos3t+\text{usw.}\\
&\quad+\beta'\sin t+\beta''\sin2t+\beta'''\sin3t+\text{usw.}
\end{aligned}
\]
\emph{identisch} sein müssen; dadurch werden die Koeffizienten $\alpha,\alpha',\beta',\alpha'',\beta''$ usw. bekannt. Ferner ist die Formel für $T$ in der hier angegebenen Gestalt so beschaffen, dass ohne Rechnung unmittelbar erkennbar ist: Wird für $t$ der Reihe nach $a,b,c,d$ usw. eingesetzt, so werden die vorgegebenen Werte $A,B,C,D$
usw. genau angenommen.

\InterpolationArticle{11}

Wäre die Anzahl der gegebenen Werte $A,B,C,D$ usw. um eins kleiner, als
wir vorausgesetzt haben, nämlich $=2m$, so würden sie nicht zur Bestimmung der $2m+1$ unbekannten Koeffizienten ausreichen, sofern nicht eine Beziehung zwischen ihnen bekannt wäre. In diesem Fall kann also der allgemeine Ausdruck für $T$ nur so angegeben werden, dass er eine unbekannte Konstante enthält. Die Überlegungen des vorhergehenden Artikels lassen sich hier nicht anwenden, da die Summe $W=0$, die aus Artikel 2 entnommen wurde, auf der Voraussetzung beruhte, dass die Anzahl der
Größen $A,B\mathinner{.\,.\,.}T$ gerade ist. Die im vorhergehenden Artikel für $T$ gefundene Formel genügt zwar im Allgemeinen den Werten $A,B$ usw., ob ihre Anzahl gerade
oder ungerade ist; bei geringer Aufmerksamkeit erkennt man aber, dass die Entwicklung der Sinusprodukte
in unserem Fall einen Ausdruck ergibt, der dem vorgegebenen nicht ähnelt, sondern die folgende Form hat
\SourceLeaf{291}
\[
\begin{aligned}
&\mathfrak{A}\cos\tfrac12t+\mathfrak{A}'\cos\tfrac32t+\mathfrak{A}''\cos\tfrac52t+\text{usw.}\\
&\quad+\mathfrak{B}\sin\tfrac12t+\mathfrak{B}'\sin\tfrac32t+\mathfrak{B}''\sin\tfrac52t+\text{usw.}
\end{aligned}
\]
Für diesen Fall muss daher ein anderer Weg eingeschlagen werden.

Man multipliziere $X$ mit $\cos(\tfrac12x+k)$. Setzt man als letzte Glieder in $X$ die Glieder $\alpha^m\cos mx+\beta^m\sin mx$ an, so ergibt sich
\[
\begin{aligned}
&\alpha\cos(\tfrac12x+k)+\tfrac12\alpha'\cos(\tfrac32x+k)\\&\qquad+\tfrac12\alpha''\cos(\tfrac52x+k)+\mathinner{.\,.\,.\,.}\\
&\hspace{2.01em}+\tfrac12\alpha^m\cos((m+\tfrac12)x+k)\\
&+\tfrac12\alpha'\cos(\tfrac12x-k)+\tfrac12\alpha''\cos(\tfrac32x-k)\\&\qquad+\tfrac12\alpha'''\cos(\tfrac52x-k)+\mathinner{.\,.\,.}\\
&\hspace{2.01em}+\tfrac12\alpha^m\cos((m-\tfrac12)x-k)\\
&\hspace{2em}+\tfrac12\beta'\sin(\tfrac32x+k)+\tfrac12\beta''\sin(\tfrac52x+k)+\mathinner{.\,.\,.}\\
&\hspace{2.01em}+\tfrac12\beta^m\sin((m+\tfrac12)x+k)\\
&+\tfrac12\beta'\sin(\tfrac12x-k)+\tfrac12\beta''\sin(\tfrac32x-k)\\&\qquad+\tfrac12\beta'''\sin(\tfrac52x-k)+\mathinner{.\,.\,.}\\
&\hspace{2.01em}+\tfrac12\beta^m\sin((m-\tfrac12)x-k)
\end{aligned}
\]
Wird nun in diesem Ausdruck für $x$ der Reihe nach $a,b,c,d\mathinner{.\,.\,.\,.}t$ eingesetzt, das
jeweilige Ergebnis mit
\[
\begin{aligned}
&\frac{1}{\sin\tfrac12(a-b)\sin\tfrac12(a-c)\sin\tfrac12(a-d)\mathinner{.\,.\,.\,.}\sin\tfrac12(a-t)}\\[1ex]
&\frac{1}{\sin\tfrac12(b-a)\sin\tfrac12(b-c)\sin\tfrac12(b-d)\mathinner{.\,.\,.\,.}\sin\tfrac12(b-t)}\\[1ex]
&\frac{1}{\sin\tfrac12(c-a)\sin\tfrac12(c-b)\sin\tfrac12(c-d)\mathinner{.\,.\,.\,.}\sin\tfrac12(c-t)}\\[1ex]
&\frac{1}{\sin\tfrac12(d-a)\sin\tfrac12(d-b)\sin\tfrac12(d-c)\mathinner{.\,.\,.\,.}\sin\tfrac12(d-t)}\\
&\text{usw.}\\[1ex]
&\frac{1}{\sin\tfrac12(t-a)\sin\tfrac12(t-b)\sin\tfrac12(t-c)\sin\tfrac12(t-d)\mathinner{.\,.\,.\,.}}
\end{aligned}
\]
multipliziert und die Summe der Produkte $=W$ gesetzt, so folgt aus dem, was wir in Artikel 2
(im ersten Fall, wo $m$ dasselbe war, was hier $2m+1$ ist) angegeben haben, leicht
\[
\begin{aligned}
W={}&\tfrac12(2i)^{2m}\alpha^m\cos(\tfrac12(a+b+c+d+\mathinner{.\,.\,.}+t)+k)\\
&+\tfrac12(2i)^{2m}\beta^m\sin(\tfrac12(a+b+c+d+\mathinner{.\,.\,.}+t)+k)
\end{aligned}
\]
Andererseits gilt
\endgroup
\begingroup
\renewcommand{\arraystretch}{1.45}
\SourceLeaf{292}
\[
\begin{aligned}
W={}&\frac{A\cos(\tfrac12a+k)}{\sin\tfrac12(a-b)\sin\tfrac12(a-c)\sin\tfrac12(a-d)\mathinner{.\,.\,.\,.}\sin\tfrac12(a-t)}\\[5pt]
&+\frac{B\cos(\tfrac12b+k)}{\sin\tfrac12(b-a)\sin\tfrac12(b-c)\sin\tfrac12(b-d)\mathinner{.\,.\,.\,.}\sin\tfrac12(b-t)}\\[5pt]
&+\frac{C\cos(\tfrac12c+k)}{\sin\tfrac12(c-a)\sin\tfrac12(c-b)\sin\tfrac12(c-d)\mathinner{.\,.\,.\,.}\sin\tfrac12(c-t)}\\[5pt]
&+\frac{D\cos(\tfrac12d+k)}{\sin\tfrac12(d-a)\sin\tfrac12(d-b)\sin\tfrac12(d-c)\mathinner{.\,.\,.\,.}\sin\tfrac12(d-t)}\\[5pt]
&+\text{usw.}\\[5pt]
&+\frac{T\cos(\tfrac12t+k)}{\sin\tfrac12(t-a)\sin\tfrac12(t-b)\sin\tfrac12(t-c)\sin\tfrac12(t-d)\mathinner{.\,.\,.\,.}}
\end{aligned}
\]
Setzen wir nun die Größe $k$, die völlig beliebig ist, $=-\tfrac12t$,
die Summe dagegen $a+b+c+d+\mathinner{.\,.\,.}=s$ (unter Ausschluss von $t$), so ergibt sich durch Multiplikation mit
$\sin\tfrac12(t-a)\sin\tfrac12(t-b)\sin\tfrac12(t-c)\sin\tfrac12(t-d)\mathinner{.\,.\,.\,.}$
\[
\begin{aligned}
T={}&\frac{\sin\tfrac12(t-b)\sin\tfrac12(t-c)\sin\tfrac12(t-d)\mathinner{.\,.\,.\,.}}{\sin\tfrac12(a-b)\sin\tfrac12(a-c)\sin\tfrac12(a-d)\mathinner{.\,.\,.\,.}}A\cos\tfrac12(t-a)\\[5pt]
&+\frac{\sin\tfrac12(t-a)\sin\tfrac12(t-c)\sin\tfrac12(t-d)\mathinner{.\,.\,.\,.}}{\sin\tfrac12(b-a)\sin\tfrac12(b-c)\sin\tfrac12(b-d)\mathinner{.\,.\,.\,.}}B\cos\tfrac12(t-b)\\[5pt]
&+\frac{\sin\tfrac12(t-a)\sin\tfrac12(t-b)\sin\tfrac12(t-d)\mathinner{.\,.\,.\,.}}{\sin\tfrac12(c-a)\sin\tfrac12(c-b)\sin\tfrac12(c-d)\mathinner{.\,.\,.\,.}}C\cos\tfrac12(t-c)\\[5pt]
&+\frac{\sin\tfrac12(t-a)\sin\tfrac12(t-b)\sin\tfrac12(t-c)\mathinner{.\,.\,.\,.}}{\sin\tfrac12(d-a)\sin\tfrac12(d-b)\sin\tfrac12(d-c)\mathinner{.\,.\,.\,.}}D\cos\tfrac12(t-d)\\[5pt]
&+\text{usw.}\\[5pt]
&+2^{2m-1}i^{2m}\sin\tfrac12(t-a)\sin\tfrac12(t-b)\sin\tfrac12(t-c)\sin\tfrac12(t-d)\mathinner{.\,.\,.}\\[5pt]
&\qquad\times\{\alpha^m\cos\tfrac12s+\beta^m\sin\tfrac12s\}
\end{aligned}
\]
\Needspace{31\baselineskip}
Zu dieser Formel fügen wir noch einige Bemerkungen hinzu

I. Es gilt $i^{2m}=+1$ oder $-1$, je nachdem, ob $m$ gerade oder ungerade ist.

II. Da die Formel ohne Einschränkung für jeden beliebigen Wert von $t$ gilt, muss sie nach
Entwicklung der Sinusprodukte in Sinus und Kosinus von Vielfachen der Bögen mit der Formel
für $X$, wenn für $x$ der Wert $t$ geschrieben wird, identisch sein, also mit
\[
\begin{aligned}
&\alpha+\alpha'\cos t+\alpha''\cos2t+\text{usw.}\\[5pt]
&\quad+\beta'\sin t+\beta''\sin2t+\text{usw.}
\end{aligned}
\]
III. Somit sind alle Koeffizienten $\alpha,\alpha',\beta',\alpha'',\beta''$ usw. durch die bekannten Größen und die einzige Unbekannte
$\alpha^m\cos\tfrac12s+\beta^m\sin\tfrac12s$ bestimmt; ist daher
zusätzlich eine Beziehung zwischen ihnen gegeben, so lassen sie sich alle leicht vollständig angeben.
\par\SourceLeaf{293}
\Needspace{14\baselineskip}
IV. Wird dagegen nur ein Ausdruck von derselben Form wie
\[
\begin{aligned}
&\alpha+\alpha'\cos t+\alpha''\cos2t+\text{usw.}\\[5pt]
&\quad+\beta'\sin t+\beta''\sin2t+\text{usw.}
\end{aligned}
\]
gesucht (ob er mit diesem identisch ist oder nicht), der die vorgegebenen Werte
$A,B,C,D$ usw. darstellt, wenn für $t$ der Reihe nach $a,b,c,d$ usw. eingesetzt wird, so wird die Größe
$\alpha^m\cos\tfrac12s+\beta^m\sin\tfrac12s$ aus einer unbekannten zu einer \emph{beliebigen} Größe, die man, wenn man will, sogar $=0$
setzen darf; tatsächlich ist unsere Formel für $T$ so beschaffen, dass man ihr diese Eigenschaft ohne
Rechnung unmittelbar ansieht.

\endgroup
